\documentclass[11pt]{article}
\usepackage[a4paper,margin=1in]{geometry}
\usepackage{amsmath,amssymb,bm}
\usepackage{physics}
\usepackage{mathtools}
\usepackage{hyperref}
\usepackage{enumitem}
\setlist{nosep}
\newcommand{\ee}{\mathrm{e}}
\newcommand{\ii}{\mathrm{i}}
\newcommand{\HLS}{H_{\mathrm{LS}}}
\newcommand{\Heff}{H_{\mathrm{eff}}}
\newcommand{\Hatom}{H_A}
\newcommand{\VR}{V_R}
\newcommand{\HR}{H_R}
\newcommand{\projG}{P}
\newcommand{\projE}{Q}
\title{Derivation of the Light-Shift Hamiltonian\\Using the Laser Rotating Frame}
\author{}
\date{}

\begin{document}
\maketitle

\section{Goal and physical regime}

We consider an atom with a set of ground states $\{\ket{g}\}$ and excited states $\{\ket{e}\}$, driven by a monochromatic classical electric field. Our goal is to derive an effective Hamiltonian that acts only within the ground-state manifold when the optical field is far enough detuned that the excited manifold is only virtually populated.

The derivation will follow the sequence
\begin{equation}
\boxed{
\text{Schr\"odinger picture}
\;\longrightarrow\;
\text{laser rotating frame}
\;\longrightarrow\;
\text{RWA}
\;\longrightarrow\;
\text{second-order effective Hamiltonian}
}
\label{eq:roadmap}
\end{equation}

This route is particularly convenient because the rotating-frame Hamiltonian becomes time independent after the rotating-wave approximation (RWA), so ordinary stationary perturbation theory can then be used.

\subsection*{Definitions}
\begin{enumerate}[label=\textbf{D\arabic*.}]
\item $\ket{g}$: a state in the ground-state manifold.
\item $\ket{e}$: a state in the excited-state manifold.
\item $E_g$: energy of $\ket{g}$ in the absence of the optical field.
\item $E_e$: energy of $\ket{e}$ in the absence of the optical field.
\item $\omega$: angular frequency of the applied optical field.
\item $\hbar$: reduced Planck constant.
\end{enumerate}

\section{Bare atomic Hamiltonian}

The atomic Hamiltonian is
\begin{equation}
\Hatom
=
\sum_g E_g\ket{g}\bra{g}
+
\sum_e E_e\ket{e}\bra{e}.
\label{eq:HA}
\end{equation}

We define projectors onto the ground and excited manifolds,
\begin{equation}
\projG=\sum_g\ket{g}\bra{g},
\qquad
\projE=\sum_e\ket{e}\bra{e},
\label{eq:projectors}
\end{equation}
with
\begin{equation}
\projG+\projE=1,
\qquad
\projG\projE=\projE\projG=0.
\label{eq:PQ}
\end{equation}

\section{Classical electric field and dipole interaction}

Write the real electric field as
\begin{equation}
\bm E(t)
=
\frac{1}{2}
\left(
\bm{\mathcal E}\,\ee^{-\ii\omega t}
+
\bm{\mathcal E}^{*}\,\ee^{+\ii\omega t}
\right),
\label{eq:Efield}
\end{equation}
where $\bm{\mathcal E}$ is the complex field amplitude. It can include the polarization vector and its complex phase.

The electric-dipole interaction is
\begin{equation}
V(t)=-\bm d\cdot \bm E(t),
\label{eq:dipoleinteraction}
\end{equation}
where $\bm d$ is the atomic electric-dipole operator.

The full Schr\"odinger-picture Hamiltonian is therefore
\begin{equation}
H_S(t)=\Hatom+V(t).
\label{eq:HS}
\end{equation}

\section{Raising and lowering parts of the dipole operator}

For optical transitions between the ground and excited manifolds, define
\begin{equation}
\bm d^{(+)}
\equiv
\projE\bm d\projG
=
\sum_{e,g}
\bm d_{eg}\ket{e}\bra{g},
\label{eq:dplus}
\end{equation}
where
\begin{equation}
\bm d_{eg}
\equiv
\bra{e}\bm d\ket{g}.
\label{eq:deg}
\end{equation}
Its Hermitian conjugate is
\begin{equation}
\bm d^{(-)}
\equiv
\projG\bm d\projE
=
\left(\bm d^{(+)}\right)^\dagger.
\label{eq:dminus}
\end{equation}
Thus
\begin{equation}
\bm d\simeq \bm d^{(+)}+\bm d^{(-)},
\label{eq:ddecomp}
\end{equation}
where direct dipole matrix elements within a single parity manifold are omitted.

Substituting Eqs.~\eqref{eq:Efield} and \eqref{eq:ddecomp} into Eq.~\eqref{eq:dipoleinteraction} gives four terms,
\begin{align}
V(t)
={}&-
\bm d^{(+)}\cdot \bm E^{(+)}(t)
-
\bm d^{(+)}\cdot \bm E^{(-)}(t)
\nonumber\\
&-
\bm d^{(-)}\cdot \bm E^{(+)}(t)
-
\bm d^{(-)}\cdot \bm E^{(-)}(t),
\label{eq:fourterms}
\end{align}
with
\begin{equation}
\bm E^{(+)}(t)=\frac{1}{2}\bm{\mathcal E}\ee^{-\ii\omega t},
\qquad
\bm E^{(-)}(t)=\frac{1}{2}\bm{\mathcal E}^{*}\ee^{+\ii\omega t}.
\label{eq:Epm}
\end{equation}

\section{Transform directly to the laser rotating frame}

We now introduce the unitary transformation
\begin{equation}
R(t)=\ee^{-\ii\omega t\projE}.
\label{eq:Rdef}
\end{equation}
This rotates only the excited manifold at the laser frequency $\omega$.

Define the rotating-frame state by
\begin{equation}
\ket{\psi_R(t)}=R^\dagger(t)\ket{\psi_S(t)}.
\label{eq:psir}
\end{equation}
The rotating-frame Hamiltonian is
\begin{equation}
\HR(t)
=
R^\dagger H_S(t)R
-
\ii\hbar R^\dagger \dot R.
\label{eq:HRgeneral}
\end{equation}

Because $\projE$ commutes with $\Hatom$,
\begin{equation}
R^\dagger \Hatom R=\Hatom.
\label{eq:HArot}
\end{equation}
Also,
\begin{equation}
\dot R=-\ii\omega\projE R,
\label{eq:Rdot}
\end{equation}
so
\begin{equation}
-\ii\hbar R^\dagger\dot R
=-\hbar\omega\projE.
\label{eq:frameenergy}
\end{equation}
Therefore the bare excited-state energies are shifted according to
\begin{equation}
E_e\longrightarrow E_e-\hbar\omega.
\label{eq:eshift}
\end{equation}

\section{How the RWA appears in the rotating frame}

The rotating transformation acts on the dipole raising and lowering operators as
\begin{equation}
R^\dagger\bm d^{(+)}R
=
\ee^{+\ii\omega t}\bm d^{(+)},
\label{eq:dplusrot}
\end{equation}
and
\begin{equation}
R^\dagger\bm d^{(-)}R
=
\ee^{-\ii\omega t}\bm d^{(-)}.
\label{eq:dminusrot}
\end{equation}

Therefore the four products in Eq.~\eqref{eq:fourterms} transform as
\begin{align}
\bm d^{(+)}\bm E^{(+)} &\longrightarrow \text{constant},
\label{eq:slow1}\\
\bm d^{(-)}\bm E^{(-)} &\longrightarrow \text{constant},
\label{eq:slow2}\\
\bm d^{(+)}\bm E^{(-)} &\longrightarrow \ee^{+2\ii\omega t},
\label{eq:fast1}\\
\bm d^{(-)}\bm E^{(+)} &\longrightarrow \ee^{-2\ii\omega t}.
\label{eq:fast2}
\end{align}

Thus the rotating-frame Hamiltonian has the structure
\begin{equation}
\HR(t)
=
H_{\mathrm{static}}
+
H_{+2\omega}\ee^{+2\ii\omega t}
+
H_{-2\omega}\ee^{-2\ii\omega t}.
\label{eq:HRstructure}
\end{equation}

The rotating-wave approximation neglects the rapidly oscillating $\ee^{\pm2\ii\omega t}$ terms. Physically, these terms average to nearly zero on the slower timescale relevant to the atomic dynamics when
\begin{equation}
|\Omega|,\ |\Delta|,\ \text{ground-state splittings}
\ll
\omega.
\label{eq:RWAcondition}
\end{equation}

\section{RWA Hamiltonian in the rotating frame}

Define the Rabi-frequency matrix element
\begin{equation}
\Omega_{eg}
\equiv
\frac{\bra{e}\bm d\cdot\bm{\mathcal E}\ket{g}}{\hbar}.
\label{eq:Rabi}
\end{equation}
With this convention, the factor $1/2$ from the real field is not absorbed into $\Omega_{eg}$.

After the RWA, the rotating-frame Hamiltonian becomes
\begin{equation}
\boxed{
\HR^{\mathrm{RWA}}
=
\Hatom
-
\hbar\omega\projE
-
\frac{\hbar}{2}
\sum_{e,g}
\left(
\Omega_{eg}\ket{e}\bra{g}
+
\Omega_{eg}^{*}\ket{g}\bra{e}
\right)
}.
\label{eq:HRRWA}
\end{equation}

This Hamiltonian is time independent.

We now write
\begin{equation}
\HR^{\mathrm{RWA}}
=
H_{0R}+\VR,
\label{eq:Hsplit}
\end{equation}
where
\begin{equation}
H_{0R}
=
\Hatom-\hbar\omega\projE,
\label{eq:H0R}
\end{equation}
and
\begin{equation}
\VR
=
-
\frac{\hbar}{2}
\sum_{e,g}
\left(
\Omega_{eg}\ket{e}\bra{g}
+
\Omega_{eg}^{*}\ket{g}\bra{e}
\right).
\label{eq:VR}
\end{equation}

\section{Detuning as a rotating-frame energy difference}

Define the optical transition frequency
\begin{equation}
\omega_{eg}
\equiv
\frac{E_e-E_g}{\hbar},
\label{eq:omegaeg}
\end{equation}
and the detuning
\begin{equation}
\boxed{
\Delta_{eg}
\equiv
\omega_{eg}-\omega
}.
\label{eq:detuning}
\end{equation}

The rotating-frame energy of $\ket{e}$ is
\begin{equation}
E_e^{(R)}=E_e-\hbar\omega.
\label{eq:Eerot}
\end{equation}
Hence
\begin{align}
E_g-E_e^{(R)}
&=
E_g-(E_e-\hbar\omega)
\nonumber\\
&=
-\hbar(\omega_{eg}-\omega)
\nonumber\\
&=
-\hbar\Delta_{eg}.
\label{eq:denominator}
\end{align}
Thus the familiar $1/\Delta_{eg}$ light-shift denominator is simply the inverse energy mismatch between the ground state and the virtually accessed excited state in the laser rotating frame.

\section{Why an effective Hamiltonian is needed}

Suppose the field is sufficiently far detuned that
\begin{equation}
|\Omega_{eg}|\ll |\Delta_{eg}|.
\label{eq:dispersivecondition}
\end{equation}
Then the excited manifold is only weakly populated. Nevertheless, the second-order virtual process
\begin{equation}
\ket{g}
\longrightarrow
\ket{e}
\longrightarrow
\ket{g'}
\label{eq:virtual}
\end{equation}
produces both diagonal energy shifts and off-diagonal couplings within the ground manifold.

Therefore we seek a ground-manifold effective Hamiltonian
\begin{equation}
\Heff^{(g)}
=
\projG H_{0R}\projG
+
H_{\mathrm{eff}}^{(2)}
+
O(\VR^3).
\label{eq:Heffg}
\end{equation}
Because $\projG\projE=0$,
\begin{equation}
\projG H_{0R}\projG
=
\projG\Hatom\projG.
\label{eq:PH0P}
\end{equation}
The second-order term is the coherent light-shift Hamiltonian,
\begin{equation}
\HLS\equiv H_{\mathrm{eff}}^{(2)}.
\label{eq:HLSdef}
\end{equation}

\section{Systematic second-order block diagonalization}

We now derive the general Hermitian second-order effective Hamiltonian rather than guessing its off-diagonal form.

Let
\begin{equation}
H=H_0+V,
\label{eq:Hgeneric}
\end{equation}
where $H_0$ is diagonal in the retained and eliminated subspaces, and $V$ couples them.

Introduce an anti-Hermitian generator $S$,
\begin{equation}
S^\dagger=-S,
\label{eq:Santi}
\end{equation}
and define
\begin{equation}
\widetilde H
=\ee^S H\ee^{-S}.
\label{eq:Htilde}
\end{equation}
Expanding to second order,
\begin{equation}
\widetilde H
=
H_0+V+[S,H_0]+[S,V]
+\frac{1}{2}[S,[S,H_0]]
+O(V^3).
\label{eq:BCH}
\end{equation}

Choose $S$ such that the first-order coupling between retained and eliminated subspaces vanishes,
\begin{equation}
V+[S,H_0]=0
\label{eq:Scancel}
\end{equation}
for the off-diagonal blocks.

For $\ket{i}$ in the retained subspace and $\ket{m}$ in the eliminated subspace,
\begin{equation}
\bra{m}[S,H_0]\ket{i}
=
(E_i-E_m)S_{mi}.
\label{eq:commmatrix}
\end{equation}
Equation~\eqref{eq:Scancel} gives
\begin{equation}
S_{mi}
=
\frac{V_{mi}}{E_m-E_i},
\qquad
S_{im}
=
\frac{V_{im}}{E_i-E_m}.
\label{eq:Smatrix}
\end{equation}

Using $[S,H_0]=-V$, Eq.~\eqref{eq:BCH} becomes
\begin{equation}
\widetilde H
=
H_0+\frac{1}{2}[S,V]+O(V^3).
\label{eq:Htilde2}
\end{equation}
Therefore
\begin{equation}
H_{\mathrm{eff}}^{(2)}
=
\frac{1}{2}\projG[S,V]\projG.
\label{eq:Heff2operator}
\end{equation}

For two retained states $\ket{i}$ and $\ket{j}$,
\begin{align}
\bra{i}H_{\mathrm{eff}}^{(2)}\ket{j}
&=
\frac{1}{2}
\sum_m
\left(
S_{im}V_{mj}-V_{im}S_{mj}
\right)
\nonumber\\
&=
\frac{1}{2}
\sum_m
V_{im}V_{mj}
\left[
\frac{1}{E_i-E_m}
+
\frac{1}{E_j-E_m}
\right].
\label{eq:generalHeff2}
\end{align}

Thus
\begin{equation}
\boxed{
\bra{i}H_{\mathrm{eff}}^{(2)}\ket{j}
=
\frac{1}{2}
\sum_m
V_{im}V_{mj}
\left[
\frac{1}{E_i-E_m}
+
\frac{1}{E_j-E_m}
\right]
}.
\label{eq:generalboxed}
\end{equation}

The symmetrization of the two denominators is essential for Hermiticity when $E_i\neq E_j$.

\section{Apply the general formula to the light shift}

Now identify
\begin{equation}
\ket{i}=\ket{g'},
\qquad
\ket{j}=\ket{g},
\qquad
\ket{m}=\ket{e},
\label{eq:identification}
\end{equation}
and use the rotating-frame unperturbed excited-state energy
\begin{equation}
E_m=E_e-\hbar\omega.
\label{eq:Emrot}
\end{equation}

From Eq.~\eqref{eq:VR},
\begin{equation}
\bra{e}\VR\ket{g}
=-\frac{\hbar}{2}\Omega_{eg},
\label{eq:Veg}
\end{equation}
and
\begin{equation}
\bra{g'}\VR\ket{e}
=-\frac{\hbar}{2}\Omega_{eg'}^*.
\label{eq:Vge}
\end{equation}
Thus
\begin{equation}
\bra{g'}\VR\ket{e}
\bra{e}\VR\ket{g}
=
\frac{\hbar^2}{4}
\Omega_{eg'}^*\Omega_{eg}.
\label{eq:Vproduct}
\end{equation}

The denominators are
\begin{equation}
E_g-(E_e-\hbar\omega)
=-\hbar\Delta_{eg},
\label{eq:deng}
\end{equation}
and
\begin{equation}
E_{g'}-(E_e-\hbar\omega)
=-\hbar\Delta_{eg'}.
\label{eq:dengp}
\end{equation}

Substituting into Eq.~\eqref{eq:generalboxed}, we obtain
\begin{equation}
\boxed{
\bra{g'}\HLS\ket{g}
=
-\frac{\hbar}{8}
\sum_e
\Omega_{eg'}^*\Omega_{eg}
\left(
\frac{1}{\Delta_{eg'}}
+
\frac{1}{\Delta_{eg}}
\right)
}.
\label{eq:HLSfinal}
\end{equation}

This is the general second-order coherent light-shift Hamiltonian in the ground manifold after the RWA.

\section{Diagonal AC Stark shift}

For $g'=g$, Eq.~\eqref{eq:HLSfinal} becomes
\begin{align}
\bra{g}\HLS\ket{g}
&=
-\frac{\hbar}{8}
\sum_e
|\Omega_{eg}|^2
\left(
\frac{1}{\Delta_{eg}}
+
\frac{1}{\Delta_{eg}}
\right)
\nonumber\\
&=
-\frac{\hbar}{4}
\sum_e
\frac{|\Omega_{eg}|^2}{\Delta_{eg}}.
\label{eq:diagonalshift}
\end{align}
Therefore
\begin{equation}
\boxed{
\delta E_g
=
-\frac{\hbar}{4}
\sum_e
\frac{|\Omega_{eg}|^2}{\Delta_{eg}}
}.
\label{eq:ACStark}
\end{equation}

With the convention $\Delta_{eg}=\omega_{eg}-\omega$, red detuning corresponds to $\Delta_{eg}>0$, so the ground-state shift is negative.

\section{Hermiticity of the off-diagonal light shift}

Taking the complex conjugate of Eq.~\eqref{eq:HLSfinal},
\begin{align}
\left[\bra{g'}\HLS\ket{g}\right]^*
&=
-\frac{\hbar}{8}
\sum_e
\Omega_{eg'}\Omega_{eg}^*
\left(
\frac{1}{\Delta_{eg'}}
+
\frac{1}{\Delta_{eg}}
\right)
\nonumber\\
&=
\bra{g}\HLS\ket{g'},
\label{eq:Hermitianproof}
\end{align}
provided the detunings are real, as they are for the coherent Hamiltonian considered here. Hence
\begin{equation}
\boxed{
\HLS^\dagger=\HLS
}.
\label{eq:HLShermitian}
\end{equation}

A frequently written unsymmetrized form,
\begin{equation}
\bra{g'}\widetilde H_{\mathrm{LS}}\ket{g}
\propto
\sum_e
\frac{\Omega_{eg'}^*\Omega_{eg}}{\Delta_{eg}},
\label{eq:wrongunsym}
\end{equation}
is not generally Hermitian when $\Delta_{eg}\neq\Delta_{eg'}$. It is only justified when the ground-state splittings are negligible compared with the optical detuning.

\section{Nearly degenerate ground manifold}

If the ground-state splittings are much smaller than the optical detuning,
\begin{equation}
|E_g-E_{g'}|\ll \hbar|\Delta|,
\label{eq:degcond}
\end{equation}
then
\begin{equation}
\Delta_{eg}\simeq\Delta_{eg'}\simeq\Delta_e.
\label{eq:degeq}
\end{equation}
Equation~\eqref{eq:HLSfinal} reduces to
\begin{equation}
\boxed{
\bra{g'}\HLS\ket{g}
\simeq
-\frac{\hbar}{4}
\sum_e
\frac{\Omega_{eg'}^*\Omega_{eg}}{\Delta_e}
}.
\label{eq:HLSdeg}
\end{equation}
This is the form commonly used before decomposing the light shift into scalar, vector, and tensor contributions.

\section{Dipole-operator form}

Using Eq.~\eqref{eq:Rabi}, Eq.~\eqref{eq:HLSfinal} can be written as
\begin{align}
\bra{g'}\HLS\ket{g}
=
-\frac{1}{8\hbar}
\sum_e
&\bra{g'}\bm d\cdot\bm{\mathcal E}^{*}\ket{e}
\bra{e}\bm d\cdot\bm{\mathcal E}\ket{g}
\nonumber\\
&\times
\left(
\frac{1}{\Delta_{eg'}}
+
\frac{1}{\Delta_{eg}}
\right).
\label{eq:HLSdipole}
\end{align}

In the nearly degenerate ground-manifold approximation,
\begin{equation}
\bra{g'}\HLS\ket{g}
\simeq
-\frac{1}{4\hbar}
\sum_e
\frac{
\bra{g'}\bm d\cdot\bm{\mathcal E}^{*}\ket{e}
\bra{e}\bm d\cdot\bm{\mathcal E}\ket{g}
}{\Delta_e}.
\label{eq:HLSdipoledeg}
\end{equation}

\section{What was neglected by the RWA?}

The discarded rotating-frame terms oscillate as $\ee^{\pm2\ii\omega t}$. If they are retained, second-order perturbation theory generates additional counter-rotating contributions with denominators of the form
\begin{equation}
\omega_{eg}+\omega,
\label{eq:counterden}
\end{equation}
instead of
\begin{equation}
\omega_{eg}-\omega=\Delta_{eg}.
\label{eq:rotden}
\end{equation}
Near resonance,
\begin{equation}
|\omega_{eg}-\omega|\ll \omega_{eg}+\omega,
\label{eq:dencompare}
\end{equation}
so the counter-rotating contribution is much smaller. The leading correction beyond the RWA is closely related to the Bloch--Siegert shift.

\section{Conditions for the derivation}

The main assumptions are:
\begin{enumerate}[label=\textbf{A\arabic*.}]
\item \textbf{Electric-dipole approximation:}
\begin{equation}
V(t)=-\bm d\cdot\bm E(t).
\end{equation}
\item \textbf{Rotating-wave approximation:}
\begin{equation}
|\Omega|,\ |\Delta|,\ \text{relevant internal splittings}\ll \omega.
\end{equation}
\item \textbf{Dispersive perturbative regime:}
\begin{equation}
|\Omega_{eg}|\ll |\Delta_{eg}|.
\end{equation}
\item \textbf{Coherent Hamiltonian treatment:} spontaneous emission is neglected in $\HLS$. If radiative decay is important, the effective dynamics should be written as a master equation containing both a Hermitian light-shift Hamiltonian and dissipative scattering terms.
\end{enumerate}

\section{Summary of the logic}

The complete derivation can be summarized as follows.

\begin{enumerate}[label=\textbf{Step \arabic*.}]
\item Start from
\begin{equation}
H_S(t)=\Hatom-\bm d\cdot\bm E(t).
\end{equation}

\item Transform directly to the laser rotating frame using
\begin{equation}
R(t)=\ee^{-\ii\omega t\projE}.
\end{equation}

\item In that frame, the resonant dipole terms become static, whereas the counter-rotating terms acquire phases $\ee^{\pm2\ii\omega t}$.

\item Apply the RWA by neglecting those rapidly oscillating terms.

\item The resulting time-independent Hamiltonian is
\begin{equation}
\HR^{\mathrm{RWA}}
=
\Hatom-\hbar\omega\projE+\VR.
\end{equation}

\item The rotating-frame energy mismatch is
\begin{equation}
E_g-(E_e-\hbar\omega)
=-\hbar\Delta_{eg}.
\end{equation}

\item Use the general second-order effective-Hamiltonian formula
\begin{equation}
\bra{i}H_{\mathrm{eff}}^{(2)}\ket{j}
=
\frac{1}{2}
\sum_m
V_{im}V_{mj}
\left[
\frac{1}{E_i-E_m}
+
\frac{1}{E_j-E_m}
\right].
\end{equation}

\item Obtain the Hermitian light-shift Hamiltonian
\begin{equation}
\boxed{
\bra{g'}\HLS\ket{g}
=
-\frac{\hbar}{8}
\sum_e
\Omega_{eg'}^*\Omega_{eg}
\left(
\frac{1}{\Delta_{eg'}}
+
\frac{1}{\Delta_{eg}}
\right)
}.
\end{equation}
\end{enumerate}

The key conceptual advantage of the rotating-frame approach is that the detuning appears immediately as a static energy difference, while the RWA is visible as the removal of terms oscillating at approximately $2\omega$. This makes the subsequent second-order perturbation theory especially direct.

\end{document}
